\begin{frame}[fragile]{Bên trong Git: blob, tree, commit và SHA}
\begin{itemize}
\item \texttt{blob} lưu nội dung file, \texttt{tree} lưu cây thư mục, \texttt{commit} trỏ tới một \texttt{tree}.
\item Mỗi \texttt{commit} gồm tác giả, thời gian, thông điệp và SHA của \texttt{commit} cha.
\item SHA-1 định danh mỗi đối tượng: đổi 1 byte thì SHA đổi hoàn toàn.
\item Ví dụ ASPX: sửa \texttt{Default.aspx} sẽ tạo \texttt{blob} mới và \texttt{commit} mới.
\item Kiểm tra đối tượng bằng lệnh \texttt{cat-file} như ví dụ dưới đây.
\end{itemize}
\begin{lstlisting}[language=bash]
git cat-file -p HEAD
git cat-file -t HEAD
git hash-object Default.aspx
\end{lstlisting}
\end{frame}

\begin{frame}{Rebase và merge}
\begin{itemize}
\item \texttt{merge} giữ nguyên lịch sử và tạo thêm một \texttt{merge commit}.
\item \texttt{rebase} viết lại lịch sử thành một đường thẳng, log sạch hơn.
\item \texttt{merge} an toàn cho nhánh public vì không đổi SHA cũ.
\item \texttt{rebase} chỉ dùng cho nhánh riêng chưa chia sẻ với ai.
\item Cấm \texttt{rebase} nhánh public: đổi SHA, gây xung đột cho cả nhóm.
\end{itemize}
\begin{tabular}{lll}
\toprule
Tiêu chí & merge & rebase \\
\midrule
Lịch sử & phân nhánh, trung thực & đường thẳng, sạch \\
Merge commit & có & không \\
Xung đột & giải một lần & giải lại từng commit \\
Dùng khi nào & nhánh public & nhánh riêng \\
\bottomrule
\end{tabular}
\end{frame}

\begin{frame}[fragile]{Cherry-pick và bisect}
\begin{itemize}
\item \texttt{cherry-pick} sao chép một \texttt{commit} sang nhánh hiện tại.
\item Dùng khi cần lấy một bản sửa lỗi mà không \texttt{merge} cả nhánh.
\item \texttt{bisect} tìm \texttt{commit} gây lỗi bằng tìm kiếm nhị phân.
\item Kết hợp \texttt{bisect run} với lệnh build để tìm lỗi tự động.
\end{itemize}
\begin{lstlisting}[language=bash]
git cherry-pick a1b2c3d
git cherry-pick --continue
\end{lstlisting}
\begin{lstlisting}[language=bash]
git bisect start
git bisect bad
git bisect good v1.0
git bisect run msbuild AspxDemo.sln
git bisect reset
\end{lstlisting}
\end{frame}

\begin{frame}[fragile]{Hooks và LFS}
\begin{itemize}
\item \texttt{Hook} là script tự chạy khi \texttt{commit}, \texttt{push} hay \texttt{merge}.
\item Ví dụ: \texttt{pre-commit} kiểm tra build, \texttt{commit-msg} kiểm tra định dạng.
\item \texttt{LFS} lưu file lớn như DLL hay ảnh ngoài repo, giữ repo nhẹ.
\item Trong repo chỉ là con trỏ text, nội dung thật nằm trên server LFS.
\end{itemize}
\begin{lstlisting}[language=bash]
ls .git/hooks
git lfs track "*.dll"
git lfs ls-files
\end{lstlisting}
\end{frame}

\begin{frame}{Quy trình nhánh: GitFlow, GitHub Flow, Trunk-Based}
\begin{itemize}
\item GitFlow nhiều nhánh, hợp với phát hành theo phiên bản lớn.
\item GitHub Flow đơn giản: một nhánh \texttt{main} cộng feature và PR.
\item Trunk-Based commit nhỏ liên tục lên \texttt{main}, đòi hỏi CI mạnh.
\item Nhóm 3 người làm ASPX 4.8.1 nên chọn GitHub Flow.
\item Mọi thay đổi đi qua PR với ít nhất một người review.
\end{itemize}
\begin{tabular}{lll}
\toprule
Quy trình & Nhánh chính & Phù hợp \\
\midrule
GitFlow & main, develop, feature & phát hành theo bản \\
GitHub Flow & main, feature & nhóm nhỏ, web \\
Trunk-Based & main & CI mạnh, nhịp nhanh \\
\bottomrule
\end{tabular}
\end{frame}

\begin{frame}{Conventional Commits}
\begin{itemize}
\item Định dạng chung: \texttt{type(scope): mô tả ngắn gọn}.
\item \texttt{feat(auth): thêm trang đăng nhập Web Forms}
\item \texttt{fix(config): sửa connection string Web config}
\item \texttt{docs(slide): bổ sung hình git flow}
\item \texttt{test(build): thêm test msbuild pipeline}
\item Lợi ích: log dễ đọc, dễ lọc và tự sinh changelog.
\end{itemize}
\end{frame}
